\chapter{Elliptic Curves as Riemann Surfaces}
\label{ch:iv31}

The Weierstrass parametrization identifies a complex torus with a smooth projective cubic curve. Analytic addition on the torus becomes the algebraic group law of an elliptic curve.

\section*{Learning goals}
After this chapter, the reader should be able to:
\begin{itemize}
\item verify that \((\wp(z),\wp'(z))\) lies on the Weierstrass cubic;
\item extend the Weierstrass parametrization across the lattice class to the point at infinity;
\item prove that the torus-to-cubic map is bijective and hence biholomorphic under the nonsingularity hypothesis;
\item identify the point at infinity as the group identity and inversion as \(z\mapsto -z\);
\item show that torus addition agrees with the chord-and-tangent group law on the cubic;
\item relate lattice scaling, the invariants \(g_2,g_3\), and the genus-one moduli viewpoint;
\end{itemize}
\section*{Conceptual roadmap}
\[
\boxed{\text{local holomorphic structure}\;\longrightarrow\;\text{integral or series representation}\;\longrightarrow\;\text{global consequence}.}
\]

\section{Weierstrass cubic}
The equation $y^2=4x^3-g_2x-g_3$ defines a nonsingular cubic when $g_2^3-27g_3^2\ne0$.

\section{Parametrization}
The map $z\mapsto[\wp(z):\wp'(z):1]$ extends across the lattice point to the point at infinity.

\section{Point at infinity}
The unique point at infinity is the identity of the cubic group law.

\section{Injectivity modulo sign and derivative}
The pair $(\wp,\wp')$ separates torus points modulo the lattice.

\section{Surjectivity}
Every point of the smooth cubic arises from a torus point.

\section{Biholomorphism}
The extended Weierstrass map is a holomorphic bijection between compact Riemann surfaces and hence a biholomorphism.

\section{Group law}
Addition of arguments corresponds to the chord-and-tangent group law on the cubic.

\section{Invariant differential}
The form $dx/y$ pulls back to $dz$ up to the normalization determined by the chosen cubic equation.

\section{Moduli viewpoint}
Changing the lattice by scaling changes $g_2,g_3$ with weights and preserves the modular invariant $j$.

\section{Genus-one synthesis}
Complex tori, smooth cubic curves, elliptic functions, and genus-one Riemann surfaces with a base point describe the same analytic geometry.

\section{Core structural results}

\begin{theorem}[Weierstrass map lies on the cubic]\label{thm:iv31-01}
For every $z\notin\Lambda$, the pair $(\wp(z),\wp'(z))$ satisfies $y^2=4x^3-g_2x-g_3$.
\begin{proof}
This is exactly the differential equation for $\wp$.
\end{proof}
\end{theorem}

\begin{theorem}[Extension at the origin]\label{thm:iv31-02}
The map $z\mapsto[\wp(z):\wp'(z):1]$ extends holomorphically at the lattice class of zero to the projective point at infinity.
\begin{proof}
Using $\wp(z)=z^{-2}+O(z^2)$ and $\wp'(z)=-2z^{-3}+O(z)$, rescale homogeneous coordinates by a suitable power of $z$ to obtain a finite nonzero projective limit equal to the unique point at infinity.
\end{proof}
\end{theorem}

\begin{theorem}[Torus-cubic biholomorphism]\label{thm:iv31-03}
If the discriminant is nonzero, the extended Weierstrass map $\mathbb C/\Lambda\to E$ is a biholomorphism onto the smooth cubic $E$.
\begin{proof}
The map is nonconstant and holomorphic between compact Riemann surfaces. The pair $\wp,\wp'$ separates points modulo the lattice, giving injectivity; compactness makes the image closed, while nonconstant holomorphic maps are open, so the image is all of the connected cubic. A holomorphic bijection between compact Riemann surfaces has holomorphic inverse.
\end{proof}
\end{theorem}

\begin{theorem}[Compatibility of group laws]\label{thm:iv31-04}
Under the Weierstrass biholomorphism, addition on $\mathbb C/\Lambda$ corresponds to the cubic chord-and-tangent group law.
\begin{proof}
The addition formula computes the coordinates of the image of $u+v$ and matches the algebraic chord construction; the identity is the lattice origin mapped to infinity and inversion is $z\mapsto-z$, which changes the sign of $\wp'$.
\end{proof}
\end{theorem}

\section{Worked examples}

\begin{example}[Weierstrass cubic diagnostic]\label{ex:iv31-worked-01}
Give a rigorous worked analysis of Weierstrass cubic. State the relevant holomorphic, contour, series, or singularity hypothesis and derive the central conclusion. The equation $y^2=4x^3-g_2x-g_3$ defines a nonsingular cubic when $g_2^3-27g_3^2\ne0$. The decisive step is to use the complex-analytic structure globally rather than reason from real-variable intuition alone.
\end{example}

\begin{example}[Parametrization diagnostic]\label{ex:iv31-worked-02}
Give a rigorous worked analysis of Parametrization. State the relevant holomorphic, contour, series, or singularity hypothesis and derive the central conclusion. The map $z\mapsto[\wp(z):\wp'(z):1]$ extends across the lattice point to the point at infinity. The decisive step is to use the complex-analytic structure globally rather than reason from real-variable intuition alone.
\end{example}

\begin{example}[Point at infinity diagnostic]\label{ex:iv31-worked-03}
Give a rigorous worked analysis of Point at infinity. State the relevant holomorphic, contour, series, or singularity hypothesis and derive the central conclusion. The unique point at infinity is the identity of the cubic group law. The decisive step is to use the complex-analytic structure globally rather than reason from real-variable intuition alone.
\end{example}

\begin{example}[Injectivity modulo sign and derivative diagnostic]\label{ex:iv31-worked-04}
Give a rigorous worked analysis of Injectivity modulo sign and derivative. State the relevant holomorphic, contour, series, or singularity hypothesis and derive the central conclusion. The pair $(\wp,\wp')$ separates torus points modulo the lattice. The decisive step is to use the complex-analytic structure globally rather than reason from real-variable intuition alone.
\end{example}

\section{Solved dossiers}
The dossiers are canonical solved problems. The provenance ledger separately records which retained corpus rules guided each topic and which dossiers were newly designed.

\begin{problem}[Weierstrass cubic diagnostic]\label{prob:iv31-dossier-01}
Give a rigorous worked analysis of Weierstrass cubic. State the relevant holomorphic, contour, series, or singularity hypothesis and derive the central conclusion.
\end{problem}
\begin{solution}
The equation $y^2=4x^3-g_2x-g_3$ defines a nonsingular cubic when $g_2^3-27g_3^2\ne0$. The decisive step is to use the complex-analytic structure globally rather than reason from real-variable intuition alone.
\end{solution}

\begin{problem}[Parametrization diagnostic]\label{prob:iv31-dossier-02}
Give a rigorous worked analysis of Parametrization. State the relevant holomorphic, contour, series, or singularity hypothesis and derive the central conclusion.
\end{problem}
\begin{solution}
The map $z\mapsto[\wp(z):\wp'(z):1]$ extends across the lattice point to the point at infinity. The decisive step is to use the complex-analytic structure globally rather than reason from real-variable intuition alone.
\end{solution}

\begin{problem}[Point at infinity diagnostic]\label{prob:iv31-dossier-03}
Give a rigorous worked analysis of Point at infinity. State the relevant holomorphic, contour, series, or singularity hypothesis and derive the central conclusion.
\end{problem}
\begin{solution}
The unique point at infinity is the identity of the cubic group law. The decisive step is to use the complex-analytic structure globally rather than reason from real-variable intuition alone.
\end{solution}

\begin{problem}[Injectivity modulo sign and derivative diagnostic]\label{prob:iv31-dossier-04}
Give a rigorous worked analysis of Injectivity modulo sign and derivative. State the relevant holomorphic, contour, series, or singularity hypothesis and derive the central conclusion.
\end{problem}
\begin{solution}
The pair $(\wp,\wp')$ separates torus points modulo the lattice. The decisive step is to use the complex-analytic structure globally rather than reason from real-variable intuition alone.
\end{solution}

\begin{problem}[Surjectivity diagnostic]\label{prob:iv31-dossier-05}
Give a rigorous worked analysis of Surjectivity. State the relevant holomorphic, contour, series, or singularity hypothesis and derive the central conclusion.
\end{problem}
\begin{solution}
Every point of the smooth cubic arises from a torus point. The decisive step is to use the complex-analytic structure globally rather than reason from real-variable intuition alone.
\end{solution}

\begin{problem}[Biholomorphism diagnostic]\label{prob:iv31-dossier-06}
Give a rigorous worked analysis of Biholomorphism. State the relevant holomorphic, contour, series, or singularity hypothesis and derive the central conclusion.
\end{problem}
\begin{solution}
The extended Weierstrass map is a holomorphic bijection between compact Riemann surfaces and hence a biholomorphism. The decisive step is to use the complex-analytic structure globally rather than reason from real-variable intuition alone.
\end{solution}

\begin{problem}[Group law diagnostic]\label{prob:iv31-dossier-07}
Give a rigorous worked analysis of Group law. State the relevant holomorphic, contour, series, or singularity hypothesis and derive the central conclusion.
\end{problem}
\begin{solution}
Addition of arguments corresponds to the chord-and-tangent group law on the cubic. The decisive step is to use the complex-analytic structure globally rather than reason from real-variable intuition alone.
\end{solution}

\begin{problem}[Weierstrass map lies on the cubic proof dossier]\label{prob:iv31-dossier-08}
Prove the following theorem-level statement and identify the essential hypothesis: For every $z\notin\Lambda$, the pair $(\wp(z),\wp'(z))$ satisfies $y^2=4x^3-g_2x-g_3$.
\end{problem}
\begin{solution}
This is exactly the differential equation for $\wp$. This identifies the mechanism that makes the complex-analytic conclusion stronger than its real-variable analogue.
\end{solution}

\begin{problem}[Extension at the origin proof dossier]\label{prob:iv31-dossier-09}
Prove the following theorem-level statement and identify the essential hypothesis: The map $z\mapsto[\wp(z):\wp'(z):1]$ extends holomorphically at the lattice class of zero to the projective point at infinity.
\end{problem}
\begin{solution}
Using $\wp(z)=z^{-2}+O(z^2)$ and $\wp'(z)=-2z^{-3}+O(z)$, rescale homogeneous coordinates by a suitable power of $z$ to obtain a finite nonzero projective limit equal to the unique point at infinity. This identifies the mechanism that makes the complex-analytic conclusion stronger than its real-variable analogue.
\end{solution}

\begin{problem}[Torus-cubic biholomorphism proof dossier]\label{prob:iv31-dossier-10}
Prove the following theorem-level statement and identify the essential hypothesis: If the discriminant is nonzero, the extended Weierstrass map $\mathbb C/\Lambda\to E$ is a biholomorphism onto the smooth cubic $E$.
\end{problem}
\begin{solution}
The map is nonconstant and holomorphic between compact Riemann surfaces. The pair $\wp,\wp'$ separates points modulo the lattice, giving injectivity; compactness makes the image closed, while nonconstant holomorphic maps are open, so the image is all of the connected cubic. A holomorphic bijection between compact Riemann surfaces has holomorphic inverse. This identifies the mechanism that makes the complex-analytic conclusion stronger than its real-variable analogue.
\end{solution}

\begin{problem}[Compatibility of group laws proof dossier]\label{prob:iv31-dossier-11}
Prove the following theorem-level statement and identify the essential hypothesis: Under the Weierstrass biholomorphism, addition on $\mathbb C/\Lambda$ corresponds to the cubic chord-and-tangent group law.
\end{problem}
\begin{solution}
The addition formula computes the coordinates of the image of $u+v$ and matches the algebraic chord construction; the identity is the lattice origin mapped to infinity and inversion is $z\mapsto-z$, which changes the sign of $\wp'$. This identifies the mechanism that makes the complex-analytic conclusion stronger than its real-variable analogue.
\end{solution}

\begin{problem}[Chapter synthesis]\label{prob:iv31-dossier-12}
Connect the definitions and principal theorems of this chapter into one reusable complex-analysis strategy.
\end{problem}
\begin{solution}
The Weierstrass parametrization identifies a complex torus with a smooth projective cubic curve. Analytic addition on the torus becomes the algebraic group law of an elliptic curve. The reusable strategy is to identify the analytic domain, choose the correct local expansion or contour representation, apply the structural theorem, and then audit singularities and topology.
\end{solution}

\section{Exercises with complete solutions}

\begin{exercise}\label{exr:iv31-01}
State and apply the central principle behind Weierstrass cubic. Give one concrete consequence.
\end{exercise}
\begin{hint}
Substitute \(x=\wp(z)\), \(y=\wp'(z)\) into the differential equation before doing any projective geometry.
\end{hint}
\begin{solution}
The equation $y^2=4x^3-g_2x-g_3$ defines a nonsingular cubic when $g_2^3-27g_3^2\ne0$. The same principle can then be reused in later contour, residue, conformal, or Riemann-surface arguments.
\end{solution}

\begin{exercise}\label{exr:iv31-02}
State and apply the central principle behind Parametrization. Give one concrete consequence.
\end{exercise}
\begin{hint}
Near the lattice origin, use \(\wp(z)=z^{-2}+O(z^2)\) and \(\wp'(z)=-2z^{-3}+O(z)\) to identify the limiting projective point.
\end{hint}
\begin{solution}
The map $z\mapsto[\wp(z):\wp'(z):1]$ extends across the lattice point to the point at infinity. The same principle can then be reused in later contour, residue, conformal, or Riemann-surface arguments.
\end{solution}

\begin{exercise}\label{exr:iv31-03}
State and apply the central principle behind Point at infinity. Give one concrete consequence.
\end{exercise}
\begin{hint}
The discriminant condition \(g_2^3-27g_3^2\ne0\) is what guarantees the cubic is smooth; check it before invoking compact Riemann-surface arguments.
\end{hint}
\begin{solution}
The unique point at infinity is the identity of the cubic group law. The same principle can then be reused in later contour, residue, conformal, or Riemann-surface arguments.
\end{solution}

\begin{exercise}\label{exr:iv31-04}
State and apply the central principle behind Injectivity modulo sign and derivative. Give one concrete consequence.
\end{exercise}
\begin{hint}
To prove injectivity, use that the pair \((\wp,\wp')\) separates points modulo the lattice, not \(\wp\) alone.
\end{hint}
\begin{solution}
The pair $(\wp,\wp')$ separates torus points modulo the lattice. The same principle can then be reused in later contour, residue, conformal, or Riemann-surface arguments.
\end{solution}

\begin{exercise}\label{exr:iv31-05}
State and apply the central principle behind Surjectivity. Give one concrete consequence.
\end{exercise}
\begin{hint}
A nonconstant holomorphic map from a compact surface has closed image; combine this with openness to obtain surjectivity onto the connected cubic.
\end{hint}
\begin{solution}
Every point of the smooth cubic arises from a torus point. The same principle can then be reused in later contour, residue, conformal, or Riemann-surface arguments.
\end{solution}

\begin{exercise}\label{exr:iv31-06}
State and apply the central principle behind Biholomorphism. Give one concrete consequence.
\end{exercise}
\begin{hint}
Once the map is a holomorphic bijection between compact Riemann surfaces, use the standard inverse-holomorphicity theorem for biholomorphism.
\end{hint}
\begin{solution}
The extended Weierstrass map is a holomorphic bijection between compact Riemann surfaces and hence a biholomorphism. The same principle can then be reused in later contour, residue, conformal, or Riemann-surface arguments.
\end{solution}

\begin{exercise}\label{exr:iv31-07}
State and apply the central principle behind Group law. Give one concrete consequence.
\end{exercise}
\begin{hint}
Under inversion \(z\mapsto-z\), \(\wp\) is unchanged and \(\wp'\) changes sign; this matches the cubic group inverse.
\end{hint}
\begin{solution}
Addition of arguments corresponds to the chord-and-tangent group law on the cubic. The same principle can then be reused in later contour, residue, conformal, or Riemann-surface arguments.
\end{solution}

\begin{exercise}\label{exr:iv31-08}
State and apply the central principle behind Invariant differential. Give one concrete consequence.
\end{exercise}
\begin{hint}
Use the addition formula to compare torus addition with the chord-and-tangent coordinates, including the point at infinity as identity.
\end{hint}
\begin{solution}
The form $dx/y$ pulls back to $dz$ up to the normalization determined by the chosen cubic equation. The same principle can then be reused in later contour, residue, conformal, or Riemann-surface arguments.
\end{solution}

\section*{Chapter summary}
The chapter now forms part of the canonical complex-analysis chain from holomorphicity through residues, global continuation, and Riemann surfaces.
